\documentclass[11pt,letterpaper]{article}

\usepackage[margin=1in]{geometry}
\usepackage{amsmath,amssymb}
\usepackage{titling}

\linespread{1.15}
\setlength{\parindent}{0pt}
\setlength{\parskip}{0.7em}

\setlength{\droptitle}{-2em}
\pretitle{\centering\Large\bfseries\setlength{\parskip}{0pt}}
\posttitle{\par\vspace{0.3em}}
\preauthor{\centering\normalsize\setlength{\parskip}{0pt}}
\postauthor{\par}
\predate{}
\postdate{}

\title{%
  Assignment 3: Derivative of a Gaussian Process\\[0.25em]
  {\normalsize\normalfont CSDS 611 --- Advanced Topics and Methods of Data Science I}%
}
\author{Andrei Cozma}
\date{}

\begin{document}
\maketitle

Let $f\sim\mathrm{GP}_{\mathbb R}(\mu,C)$, where $\mu$ is continuously
differentiable and
\[
  C(x,x')=\sigma^2\exp\!\left(-\frac{(x-x')^2}{2\ell^2}\right),
  \qquad x,x'\in\mathbb R,\quad \sigma^2\geq0,\quad\ell>0.
\]
I interpret $f'$ as a mean-square derivative: the difference quotients
converge with expected squared error tending to zero.
For $h\ne0$, define the finite difference
\[
  D_h(x)=\frac{f(x+h)-f(x)}h.
\]
To prove that $f'$ is a Gaussian process, its values at every finite
collection of inputs must be jointly Gaussian.
For any $N$ and inputs $x_1,\ldots,x_N$, the vector
$(D_h(x_1),\ldots,D_h(x_N))^t$ is a linear transformation of the
jointly Gaussian values of $f$ at $x_i$ and $x_i+h$.
It is therefore jointly Gaussian.

Linearity of expectation gives
\[
  \mu_h(x):=\mathbb E[D_h(x)]
  =\frac{\mu(x+h)-\mu(x)}h
  \longrightarrow\mu'(x).
\]
Using two step sizes lets us compare different approximations and establish
that a derivative limit exists. For $h,r\ne0$, expanding the covariance gives
\begin{align*}
  R_{h,r}(x,x')
  &:=\operatorname{Cov}(D_h(x),D_r(x'))\\
  &=\frac{C(x+h,x'+r)-C(x+h,x')-C(x,x'+r)+C(x,x')}{hr}.
\end{align*}
Write $C_{\partial}(x,x'):=\partial^2 C(x,x')/(\partial x\,\partial x')$.
Since these mixed derivatives are continuous, the fundamental theorem of
calculus gives
\[
  R_{h,r}(x,x')
  =\int_0^1\!\int_0^1 C_{\partial}(x+th,x'+sr)\,ds\,dt
  \longrightarrow C_{\partial}(x,x')\qquad(h,r\to0).
\]
The integral averages the continuous mixed derivative at points approaching
$(x,x')$, so it converges to $C_{\partial}(x,x')$.
For each fixed $x$, it follows that
\begin{align*}
  \mathbb E\bigl[|D_h(x)-D_r(x)|^2\bigr]
  &=R_{h,h}(x,x)+R_{r,r}(x,x)-2R_{h,r}(x,x)\\
  &\quad+\bigl(\mu_h(x)-\mu_r(x)\bigr)^2
  \longrightarrow0\qquad(h,r\to0).
\end{align*}
Each covariance tends to $C_{\partial}(x,x)$, so the first three terms
cancel in the limit. Both means tend to $\mu'(x)$, so the final term
also tends to zero.
Thus the difference quotients satisfy the Cauchy criterion for convergence
in mean square. Their limit, denoted by $f'(x)$, satisfies
\[
  \mathbb E\bigl[|D_h(x)-f'(x)|^2\bigr]\longrightarrow0.
\]
By the Cauchy--Schwarz inequality, mean-square convergence implies convergence
of means and covariances. Therefore,
\[
  \mathbb E[f'(x)]=\mu'(x),
  \qquad
  \operatorname{Cov}(f'(x),f'(x'))=C_{\partial}(x,x').
\]
For any finite collection of inputs, the Gaussian vectors of finite
differences converge in mean square to the vector of derivative values.
Their mean vectors and covariance matrices converge to those just
calculated. Gaussian distributions with convergent means and covariance
matrices have a Gaussian limiting distribution, so these derivative
values are jointly Gaussian. Since the inputs and $N$ were arbitrary,
$f'$ is a Gaussian process.

Differentiating the squared exponential kernel gives
\begin{align*}
  \frac{\partial C(x,x')}{\partial x}
  &=-\frac{x-x'}{\ell^2}C(x,x'),\\
  C_{\partial}(x,x')
  &=\frac{\partial}{\partial x'}
    \left[-\frac{x-x'}{\ell^2}C(x,x')\right]\\
  &=\left(\frac1{\ell^2}-\frac{(x-x')^2}{\ell^4}\right)C(x,x').
\end{align*}
The kernel is symmetric, and the derivative variance is
$C_{\partial}(x,x)=\sigma^2/\ell^2$.
Positive semidefiniteness follows because $C_{\partial}$ is the covariance
of the derivative process.

Therefore,
\[
  \boxed{\begin{gathered}
    f'\sim\mathrm{GP}_{\mathbb R}(\mu',C_{\partial}),\\[0.4em]
    C_{\partial}(x,x')=
    \sigma^2\left(\frac1{\ell^2}-\frac{(x-x')^2}{\ell^4}\right)
    \exp\!\left(-\frac{(x-x')^2}{2\ell^2}\right).
  \end{gathered}}
\]

\end{document}
